% =============================================================================
%  lec18.tex —— 第 18 讲 Dold–Kan 对应（底本 PDF 第 139–145 页）
% =============================================================================
\lecture{18}{Dold–Kan 对应}
\begin{center}\small\color{black!55}（底本 PDF 第 139–145 页）\end{center}

本讲讨论 Dold–Kan 对应及其各种推广。我们采用同调约定，因为这在关于 DK 对应的文献中是标准做法。

\yizhu{“同调约定”指链复形的次数用下标书写（$C_{n}$、$N_{n}$），微分降低次数；这与本讲义前面各讲（如第 2、3 讲）所用的上同调约定（$\Ch^{\le 0}$、上标）相反。原文明确声明采用同调约定，故本讲一律保留下标写法。\textbf{[来源：B]}}

\section{经典 Dold–Kan 对应}

\nitem\label{L18-1.1} 设 $A$ 是 Abel 范畴。经典的 Dold–Kan 对应是下述两者之间的等价：
\begin{itemize}[leftmargin=2.2em,itemsep=0.1em]
\item $A$ 中非负分次链复形构成的普通范畴 $\Ch(A)_{\ge 0}$；
\item $A$ 中单纯对象构成的普通范畴 $A_{\Delta}$。
\end{itemize}
这个等价可以沿两个方向显式地构造出来。

\begin{construction}\label{L18-1.2}
设 $A_{\bullet}\in A_{\Delta}\defeq\Fun(\Delta^{\op},A)$ 是加性范畴 $A$ 中的单纯对象。对任意 $i\in[n]$，考虑面态射 $d_{i}\colon A_{n}\to A_{n-1}$，它对应唯一的单射 $\delta^{i}_{n}\colon[n-1]\to[n]$，该单射的像不含 $i$。定义
\[
C_{*}(A)\defeq[\cdots\to A_{2}\to A_{1}\to A_{0}\to 0\to\cdots]
\]
其中
\begin{itemize}[leftmargin=2.2em,itemsep=0.1em]
\item 对象 $A_{n}$ 放在同调次数 $n$（$=$ 上同调次数 $-n$）处；
\item 微分 $A_{n}\to A_{n-1}$ 是交错和 $\sum_{i=0}^{n}(-1)^{i}d_{i}$。
\end{itemize}
我们称之为与 $A_{\bullet}$ 相伴的\textbf{非规范化链复形}。
\end{construction}

\begin{exercise}\label{L18-1.3}
证明 $C_{*}(A)$ 确实是复形。
\end{exercise}

\nitem\label{L18-1.4} 由上述练习，我们有 $C_{*}(A)\in\Ch(A)_{\ge 0}$。然而，所得到的函子 $A_{\Delta}\to\Ch(A)_{\ge 0}$ 并不是等价。

\begin{exercise}\label{L18-1.5}
取 $A\defeq\Ab$，考虑常值单纯 Abel 群 $\Z$。证明
\[
C_{*}(\Z)\simeq[\cdots\xrightarrow{0}\Z\xrightarrow{1}\Z\xrightarrow{0}\Z\to 0\to\cdots].
\]
并推出函子 $C_{*}$ 连全忠实都不是。
\end{exercise}

\nitem\label{L18-1.6} 在上述例子中，注意：
\begin{itemize}[leftmargin=2.2em,itemsep=0.1em]
\item 复形 $[\cdots\to 0\to\Z\to 0\to\cdots]$ 具有正确的自同态集合；
\item 上述复形拟同构于 $C_{*}(\Z)$；
\item 当 $n>1$ 时，$\Z$ 中的 $n$-单形都是退化的；
\end{itemize}
这表明 $C_{*}$ 不是等价，原因在于退化单形。摆脱退化单形有两种办法。

\begin{construction}\label{L18-1.7}
设 $A_{\bullet}\in A_{\Delta}$ 是 Abel 范畴 $A$ 中的单纯对象。定义
\[
D_{n}(A)\defeq\im\Big(\bigoplus_{i=0}^{n-1}A_{n-1}\to A_{n}\Big),
\]
其中对 $0\le i\le n-1$，态射 $A_{n-1}\to A_{n}$ 是第 $i$ 个退化态射。
\end{construction}

\begin{exercise}\label{L18-1.8}
证明 $D_{*}(A)$ 是 $C_{*}(A)$ 的子复形。特别地，我们可以构造商复形 $C_{*}(A)/D_{*}(A)$。
\end{exercise}

\begin{construction}\label{L18-1.9}
设 $A_{\bullet}\in A_{\Delta}$ 是 Abel 范畴 $A$ 中的单纯对象。定义
\[
N_{*}(A)\defeq[\cdots\to N_{2}(A)\to N_{1}(A)\to N_{0}(A)\to 0\to\cdots]
\]
其中
\begin{itemize}[leftmargin=2.2em,itemsep=0.1em]
\item 对象 $N_{n}(A)$ 是 $d_{i}$（$1\le i\le n$）的共同核；
\item 微分由 $d_{0}$ 给出。
\end{itemize}
我们称之为与 $A_{\bullet}$ 相伴的\textbf{规范化链复形}，或称 \textbf{Moore 复形}。
\end{construction}

\begin{exercise}\label{L18-1.10}
证明 $d_{0}$ 确实把 $N_{n}(A)$ 映入 $N_{n-1}(A)$，且 $N_{*}(A)$ 是 $C_{*}(A)$ 的子复形。
\end{exercise}

\begin{exercise}\label{L18-1.11}
证明 $C_{*}(A)\simeq N_{*}(A)\oplus D_{*}(A)$。
\end{exercise}

\nitem\label{L18-1.12} 特别地，复合
\[
N_{*}(A)\to C_{*}(A)\to C_{*}(A)/D_{*}(A)
\]
是 $\Ch(A)^{\le 0}$ 中的同构。下述结果由 Eilenberg–MacLane 证明。

\yizhu{原文此处把 $\Ch(A)$ 的次数指标写成上标 $^{\le 0}$，而本讲其余各处（1.4、1.15、1.19）均写作下标 $\ge 0$，附录 B.1 又写作下标 $\le 0$。原文记号并不统一，译本一律照录、不作改动，读者引用时请以语境为准。\textbf{[来源：I]}}

\begin{proposition}[HA.1.2.3.17]\label{L18-1.13}
链态射 $N_{*}(A)\to C_{*}(A)$ 是拟同构。
\end{proposition}

\begin{exercise}\label{L18-1.14}
事实上，$N_{*}(A)\to C_{*}(A)$ 甚至是链同伦等价。见 [GJ09, Theorem III.2.4]。
\end{exercise}

\nitem\label{L18-1.15} 现在我们来构造函子 $N_{*}\colon A_{\Delta}\to\Ch(A)_{\ge 0}$ 的逆。

\begin{construction}\label{L18-1.16}
设 $A$ 是加性范畴，$A_{*}\in\Ch(A)_{\ge 0}$。如下定义 $\mathrm{DK}_{\bullet}(A)\in A_{\Delta}$。
\begin{itemize}[leftmargin=2.2em,itemsep=0.1em]
\item 对每个 $n\ge 0$，$n$-单形构成的对象是
\[
\mathrm{DK}_{n}(A)\defeq\bigoplus_{\alpha\colon[n]\onto[k]}A_{k},
\]
其中直和按满射 $\alpha\colon[n]\to[k]$ 取指标。
\item 对 $\Delta$ 中的 $[n']\to[n]$，连接态射 $\mathrm{DK}_{n}(A)\to\mathrm{DK}_{n'}(A)$ 由映射
\[
\bigoplus_{\alpha\colon[n]\onto[k]}A_{k}\to\bigoplus_{\alpha\colon[n']\onto[k']}A_{k'}
\]
给出，其中
\begin{itemize}[leftmargin=2.2em,itemsep=0.1em,label={--}]
\item 每个交换方形
\begin{center}
\begin{tikzcd}[column sep=3.4em,row sep=1.8em]
{[n']}\arrow[r]\arrow[d] & {[n]}\arrow[d]\\
{[k]}\arrow[r] & {[k]}
\end{tikzcd}
\end{center}
贡献 $A_{k}\xrightarrow{\id}A_{k}$；
\item 每个交换方形
\begin{center}
\begin{tikzcd}[column sep=3.4em,row sep=1.8em]
{[n']}\arrow[r]\arrow[d] & {[n]}\arrow[d]\\
{[k-1]}\arrow[r,"\delta^{0}_{k}"] & {[k]}
\end{tikzcd}
\end{center}
贡献 $A_{k}\xrightarrow{d}A_{k-1}$。
\end{itemize}
\end{itemize}
可以验证 $\mathrm{DK}_{\bullet}(A)$ 确实是 $A$ 中的单纯对象。
\end{construction}

\begin{exercise}\label{L18-1.17}
取 $A\defeq\Ab$，证明 $\mathrm{DK}_{\bullet}(\Z[n])\simeq\Z\Delta^{n}/\Z\partial\Delta^{n}$，其中对单纯集 $S$，$\Z S$ 是自由单纯 Abel 群。
\end{exercise}

\nitem\label{L18-1.18} 下述结果即所谓的经典 Dold–Kan 对应。

\begin{theorem}[HA.1.2.3.7, 1.2.3.12]\label{L18-1.19}
设 $A$ 是加性范畴。
\begin{enumerate}[label=(\arabic*),leftmargin=2.4em,itemsep=0.1em]
\item 函子 $\mathrm{DK}\colon\Ch(A)_{\ge 0}\to A_{\Delta}$ 是全忠实的。
\item 若 $A$ 是幂等完备的，则 $\mathrm{DK}$ 是等价。特别地，当 $A$ 是 Abel 范畴时，$\mathrm{DK}$ 是等价。
\item 若 $A$ 是 Abel 范畴，则 $N\colon A_{\Delta}\to\Ch(A)_{\ge 0}$ 是 $\mathrm{DK}$ 的逆。
\end{enumerate}
\end{theorem}

\begin{exercise}\label{L18-1.20}
当 $A=\Ab$ 时，证明函子
\[
\Ch(\Ab)_{\ge 0}\to\Ab_{\Delta},\qquad A\mapsto\Hom(N(\Z\Delta^{\bullet}),A)
\]
是 $N$ 的右伴随。推出这个函子等价于 $\mathrm{DK}$。你能直接证明这个等价吗？
\end{exercise}

\section{模型范畴的 Dold–Kan 对应}

\nitem\label{L18-2.1} 回忆 $\Ch(\Ab)_{\ge 0}$ 与 $\Ab_{\Delta}$ 都带有经典模型结构，并且 Dold–Kan 等价以最强的意义与这些模型结构相容。本节的结论可在 [Qui67, §II.4] 中找到。

\begin{proposition}\label{L18-2.2}
$\Ch(\Ab)_{\ge 0}$ 上有一个模型结构，称为\textbf{投射模型结构}，使得：
\begin{itemize}[leftmargin=2.2em,itemsep=0.1em]
\item[(W)] 弱等价是拟同构；
\item[(C)] 上纤维化是逐次单射，且其核为投射对象；
\item[(F)] 纤维化是逐次满射。
\end{itemize}
\end{proposition}

\begin{proposition}\label{L18-2.3}
$\Ab_{\Delta}$ 上有一个模型结构，称为\textbf{经典模型结构}，使得：
\begin{itemize}[leftmargin=2.2em,itemsep=0.1em]
\item[(W)] 弱等价是底单纯集之间的弱同伦等价；
\item[(F)] 纤维化是底单纯集之间的 Kan 纤维化。
\end{itemize}
\end{proposition}

\begin{theorem}[Quillen]\label{L18-2.4}
取 $A\defeq\Ab$。等价 $\mathrm{DK}$ 与 $N$ 都保持并检测类 (W)、(C) 与 (F)。
\end{theorem}

\begin{exercise}\label{L18-2.5}
证明 $\Ab_{\Delta}$ 中任何对象都是纤维的，即其底单纯集是 Kan 复形。你能直接证明这一点吗？
\end{exercise}

\begin{exercise}\label{L18-2.6}
证明 $\Ch(\Ab)_{\ge 0}$ 一侧的同调群 $H_{*}(-)$ 对应于 $\Ab_{\Delta}$ 一侧的单纯同伦群 $\pi_{*}(-)$。
\end{exercise}

\section{连通链复形与空间}

\nitem\label{L18-3.1} 由 Quillen 定理，函子
\[
\Ch(\Ab)_{\ge 0}\to\Ab_{\Delta}\to(\Set_{\Delta})_{\ast}
\]
保持弱等价。因此我们得到它们的 $\infty$-范畴局部化之间的函子
\[
\Ch(\Ab)_{\ge 0}[W^{-1}]\to(\Set_{\Delta})_{\ast}[W^{-1}].
\]

\begin{exercise}\label{L18-3.2}
我们有 $\Ch(\Ab)_{\ge 0}[W^{-1}]\simeq\mathrm{D}(\Ab)_{\ge 0}$。
\end{exercise}

\begin{exercise}\label{L18-3.3}
我们有 $(\Set_{\Delta})_{\ast}[W^{-1}]\simeq\Spc_{\ast}$。
\end{exercise}

\nitem\label{L18-3.4} 我们滥用记号，把所得到的函子仍记作
\[
\mathrm{DK}\colon\mathrm{D}(\Ab)_{\ge 0}\to\Spc_{\ast}.
\]

\begin{warning}\label{L18-3.5}
函子 $\mathrm{DK}$ 不是等价，因为上述构造用到了遗忘函子 $\Ab_{\Delta}\to(\Set_{\Delta})_{\ast}$。
\end{warning}

\begin{exercise}\label{L18-3.6}
函子 $\mathrm{DK}$ 有左伴随。提示：构造 Quillen 伴随
\[
(\Set_{\Delta})_{\ast}\rightleftarrows\Ab_{\Delta}.
\]
\end{exercise}

\nitem\label{L18-3.7} 特别地，我们可以取谱对象，得到函子
\[
\Sptr(\mathrm{D}(\Ab)_{\ge 0})\to\Sptr(\Spc_{\ast}).
\]

\begin{exercise}\label{L18-3.8}
我们有 $\Sptr(\mathrm{D}(\Ab)_{\ge 0})\simeq\mathrm{D}(\Ab)$。
\end{exercise}

\nitem\label{L18-3.9} 我们滥用记号，把所得到的函子仍记作
\[
\mathrm{DK}\colon\mathrm{D}(\Ab)\to\Sptr.
\]

\begin{exercise}\label{L18-3.10}
将上述函子与第 17 讲练习 3.14 中的函子等同起来。
\end{exercise}

\section{稳定 $\infty$-范畴的 Dold–Kan 对应}

\nitem\label{L18-4.1} 对稳定 $\infty$-范畴 $\cC$，Lurie 的 Dold–Kan 对应给出下述两者之间的典范等价：
\begin{itemize}[leftmargin=2.2em,itemsep=0.1em]
\item $\cC$ 中 $\Z_{\ge 0}$-滤过对象构成的稳定 $\infty$-范畴 $\Fun(\Z_{\ge 0},\cC)$；
\item $\cC$ 中单纯对象构成的稳定 $\infty$-范畴 $\cC_{\Delta}\defeq\Fun(\Delta^{\op},\cC)$。
\end{itemize}

\begin{theorem}[HA.1.2.4.1]\label{L18-4.2}
设 $\cC$ 是稳定 $\infty$-范畴。则存在典范等价
\[
\Fun(\Z_{\ge 0},\cC)\simeq\cC_{\Delta}.
\]
\end{theorem}

\begin{remark}\label{L18-4.3}
在本注中，我们解释上述定理的内容，以及它与经典 Dold–Kan 对应的关系。

首先，对 $\cC$ 中任意滤过对象 $D(0)\to D(1)\to D(2)\to\cdots$，考虑移位分次对象
\[
C_{k}\defeq\operatorname{Cofib}(D(k-1)\to D(k))[-k],
\]
其中对 $k=0$ 取 $D(-1)\defeq 0$。对每个 $k$，态射 $D(k-1)\to D(k)\to D(k+1)$ 诱导纤维—余纤维序列
\[
C_{k}[k]\to\operatorname{Cofib}(D(k-1)\to D(k+1))\to C_{k+1}[k+1],
\]
它对应于一个态射 $C_{k+1}\to C_{k}$。

可以验证（练习！），在同伦范畴 $\mathrm{h}\cC$ 中，复合 $C_{k+1}\to C_{k}\to C_{k-1}$ 为零，于是我们得到加性范畴 $\mathrm{h}\cC$ 中的链复形
\[
\cdots\to C_{2}\to C_{1}\to C_{0}\to 0\to\cdots.
\]
应用经典函子 $\mathrm{DK}$，我们得到 $\mathrm{h}\cC$ 中的一个单纯对象。

另一方面，对 $\cC$ 中任意单纯对象 $A_{\bullet}$，函子 $\cC\to\mathrm{h}\cC$ 诱导 $\mathrm{h}\cC$ 中的一个单纯对象。现在，定理~\ref{L18-4.2} 的确切表述是：若 $D(-)$ 在所声称的等价下对应于 $A_{\bullet}$，则所得到的 $\mathrm{h}\cC$ 中的两个单纯对象被典范地等同。
\end{remark}

\begin{remark}\label{L18-4.4}
事实上，第 $n$ 项 $D(n)$ 可以计算为单纯对象 $A_{\bullet}$ 的 $k$-骨架的余极限。然而，由于 $\cC$ 不是普通范畴，我们需要小心单纯对象的 $k$-骨架的含义。尽管如此，这说明滤过对象的余极限应当同构于整个单纯对象 $A_{\bullet}$ 的余极限。换言之，
\[
\colim_{n\in\Z_{\ge 0}}D(n)\simeq\colim_{\Delta^{\op}}A_{\bullet}.
\]
这确实是对的，详见 HA.1.2.4。
\end{remark}

\begin{remark}\label{L18-4.5}
设 $\cC$ 配有 t-结构。对任意 $\Z$-滤过对象 $D(-)$，可以在 Abel 范畴 $\cC^{\heartsuit}$ 中构造一个谱序列，其 $E_{1}$-页由
\[
E^{p,q}_{1}=H_{p+q}\operatorname{Cofib}(D(p)\to D(p+1))
\]
给出。在一定条件下，该谱序列收敛到
\[
H_{n}\Big(\colim_{m}D(m)\Big).
\]
注意，这提供了计算 $\cC$ 中单纯对象的几何实现的同调的一种策略。详见 HA.1.2.3。
\end{remark}

\appendixsection{Moore 复形与单纯同伦群}

\nitem\label{L18-A.1} 对带基点的单纯集 $X$，可以构造一个单纯群 $G(X)$，称为 \textbf{Kan 环路群构造}，它模拟拓扑空间的环路构造。我们有
\[
\pi_{n}(X)\simeq H_{n-1}(NG(X)),
\]
它通过 Kan 环路群 $G(X)$ 的 Moore 复形的同调群来计算 $X$ 的单纯同伦群。注意，构造 $N_{*}(-)$ 对非 Abel 单纯群也有意义。

\nitem\label{L18-A.2} \textbf{建议阅读}\quad [Kan58]（原始文献）、[GJ09, §V.7]（较近的文献）。

\appendixsection{幺半 Dold–Kan 对应}

\nitem\label{L18-B.1} 考虑 $\Ab$ 上由 Abel 群张量积给出的对称幺半结构。它诱导下述结构：
\begin{itemize}[leftmargin=2.2em,itemsep=0.1em]
\item $\Ch(\Ab)_{\le 0}$ 上由分次张量积给出的对称幺半结构；
\item $\Ab_{\Delta}$ 上由逐点张量积给出的对称幺半结构。
\end{itemize}
Dold–Kan 对应并不严格与这些结构交换。下面列出关于函子 $N\colon\Ab_{\Delta}\to\Ch(A)_{\ge 0}$ 已知的事实。
\begin{enumerate}[label=(\roman*),leftmargin=2.4em,itemsep=0.1em]
\item 函子 $N$ 具有对称松弛幺半结构：对任意 $A,B\in\Ab_{\Delta}$，存在一个态射，称为 \textbf{Eilenberg–Zilber 映射}，
\[
\mathrm{EZ}_{A,B}\colon N(A)\otimes N(B)\to N(A\otimes B),
\]
它与交换性约束相容。
\item 函子 $N$ 具有余松弛幺半结构：对任意 $A,B\in\Ab_{\Delta}$，存在一个态射，称为 \textbf{Alexander–Whitney 映射}，
\[
\mathrm{AW}_{A,B}\colon N(A\otimes B)\to N(A)\otimes N(B),
\]
它与交换性约束\emph{不}相容。
\item 复合 $\mathrm{AW}_{A,B}\circ\mathrm{EZ}_{A,B}$ 等于恒等态射，而 $\mathrm{EZ}_{A,B}\circ\mathrm{AW}_{A,B}$ 只是链同伦等价。
\end{enumerate}

\nitem\label{L18-B.2} 注意，(ii) 说明逆函子 $\mathrm{DK}$ 具有自然的松弛幺半结构。由此
\[
\Ch(\Ab)\xrightarrow{\tau_{\ge 0}}\Ch(\Ab)_{\ge 0}\to\Ab_{\Delta}\to\Kan
\]
也具有自然的松弛幺半结构。特别地，任何 dg-范畴 $\cC$ 也以 Kan 充实。设 $\cC'$ 是对应的 Kan 充实范畴。我们有典范等价（HA.1.3.1.17）
\[
N(\cC')\simeq N_{\mathrm{dg}}(\cC).
\]

\begin{exercise}\label{L18-B.3}
利用上述等价推出第 17 讲命题 2.16。
\end{exercise}

\nitem\label{L18-B.4} Dold–Kan 对应不能成为幺半伴随，这使得比较下述理论并非平凡：
\begin{itemize}[leftmargin=2.2em,itemsep=0.1em]
\item 连通 dg-代数；
\item 单纯结合代数。
\end{itemize}
尽管仍有可能找到上述对象的模型范畴之间的 Quillen 等价，但当我们考虑连通交换 dg-代数与单纯交换代数时，情况就微妙得多。事实上，在正特征下，这两种理论并不等价，而且前者上没有合理的模型结构。此外，$\mathbb{F}_{p}$ 上的单纯交换代数理论与 $E_{\infty}$-$H\mathbb{F}_{p}$-谱的理论也不相同。这导致导出代数几何与谱代数几何之间的二分。更多信息见这个 mathoverflow 问题（尤其是 Lurie 与 May 的回答）。

\nitem\label{L18-B.5} \textbf{建议阅读}\quad 这个 nLab 页面中的参考文献。

\subsection*{参考文献}
\par\noindent\hangindent=5.2em\hangafter=1 [GJ09]\hspace{0.6em}Paul G Goerss and John F Jardine. \emph{Simplicial homotopy theory}. Springer Science \& Business Media, 2009.
\par\noindent\hangindent=5.2em\hangafter=1 [Kan58]\hspace{0.6em}Daniel M Kan. A combinatorial definition of homotopy groups. \emph{Annals of Mathematics}, 67(2):282–312, 1958.
\par\noindent\hangindent=5.2em\hangafter=1 [Qui67]\hspace{0.6em}Daniel G Quillen. \emph{Homotopical algebra}. Lecture Notes in Mathematics, 1967.
